\documentclass[10pt]{article}
\usepackage[a4paper,margin=22mm]{geometry}
\usepackage{amsmath,amssymb,booktabs,xcolor}
\setlength{\parskip}{0.35em}\setlength{\parindent}{0pt}
\makeatletter\def\eqref#1{\textcolor{blue}{(\detokenize{#1})}}\makeatother
\begin{document}
% ---------------------------------------------------------------------
% Drop-in continuation of Sec. VI-B (ground-effect-induced bias) for the
% IEEE Access manuscript.  Place after the LPV affine fit of the channel
% gains (Table 2 / Fig. 15); it takes the coefficients from there and
% adds what the excitation data say about them.
%
% Every number is produced by a script in analysis/ and is named with it
% in a comment, so each can be regenerated:
%   the gap and its channel split   analysis/mcrit_prediction.py
%   the attribution fits            analysis/static_attribution.py
%   the geometry checks             analysis/contact_lever_check.py
%   the three null tests            analysis/pivot_height_scan.py
%                                   analysis/arc_residual_test.py
%   the offset variants             analysis/offset_ge_variants.py
%   the internal consistency        analysis/offset_pivot_based.py
%
% Cross-references assumed to exist in the host document; rename before
% compiling into the manuscript:
%   \eqref{eq:mcrit_roll}   Eq. (7)    directional roll thresholds
%   \eqref{eq:mcrit_pitch}  Eq. (14)   directional pitch thresholds
%   \eqref{eq:weight}       Eq. (34)   weight from the moments
%   \eqref{eq:offset}       Eq. (35)   offset per direction
%   \eqref{eq:pivotfree}    Eq. (36)   pivot-free average
%   \eqref{eq:parallelaxis} Eq. (82)   parallel-axis identity
% ---------------------------------------------------------------------

\section*{What the excitation data say about the bias}

The coefficients above are predictions. The excitation runs test them
without any quantity being fitted, through the gap that the rigid
balance leaves.

\paragraph{The balance with the ground-effect channels restored}
At the onset the far contact unloads, and the moments about the loaded
contact balance. Ground effect acts on the two terms the rotors
produce and on neither of the others: the collective thrust is
augmented near the ground, and the commanded moment is delivered with a
gain. Gravity is not,
\begin{equation}
  \underbrace{(1+\eta_M^{GE})\,M_{\rm cmd}}_{\text{commanded}}
  \;+\;
  \underbrace{(1+\eta_f^{GE})\,f\,l_p}_{\text{thrust}}
  \;=\;
  \underbrace{W l_p + S_{\rm off}}_{\text{gravity}} ,
  \label{eq:gebalance}
\end{equation}
with $l_p$ the contact arm and $S_{\rm off}$ the offset moment. Solving
for the command gives the predicted threshold,
\begin{equation}
  M_{\rm cmd} =
  \frac{\mathrm{sgn}\,\bigl[W - (1+\eta_f^{GE})f\bigr]\,l_p
        + S_{\rm off}}{1+\eta_M^{GE}} ,
  \label{eq:gepred}
\end{equation}
which reduces to \eqref{eq:mcrit_roll} and \eqref{eq:mcrit_pitch} when
both channels vanish. Note the order in \eqref{eq:gepred}: the thrust
is corrected \emph{before} the tipping requirement is formed, because
what remains on the contact is $W - (1+\eta_f^{GE})f$ and not $W - f$.

\paragraph{The gap, and which channel supplies it}
Writing the no-ground-effect threshold as $T_{\rm rigid}$ and the
detected onset moment as $M_{\rm det}$,
\begin{equation}
  \Delta M \;\triangleq\; T_{\rm rigid} - M_{\rm det}
  \label{eq:gap}
\end{equation}
is the moment the ground supplies but the rigid balance omits. Over
the twenty direction groups of the five configurations the vehicle tips
at a \emph{smaller} commanded moment than the rigid balance requires in
every one of them, by $204$~mN\,m in the median --- about a fifth of
the threshold itself.

The interference model of \eqref{eq:ge_moment} accounts for $165$ of
that, or $81\%$, with no tuned constant. That figure should be split,
because it bundles two claims of very different standing. Linearising
\eqref{eq:gepred}, the thrust channel contributes $\eta_f^{GE}fl_p$ and
the moment channel $\eta_M^{GE}M$:
\begin{center}\small
\begin{tabular}{@{}lrrr@{}}
\hline
 & median & mean & share of $\Delta M$ \\
\hline
gap $\Delta M$                           & $204$ & $202$ & --- \\
thrust channel $\eta_f^{GE}fl_p$         & $117$ & $115$ & $57\%$ \\
moment channel $\eta_M^{GE}M$            & $43$  & $47$  & $21\%$ \\
residual                                 & $49$  & $40$  & $22\%$ \\
\hline
\end{tabular}
\end{center}
in mN\,m. Nearly three-quarters of the modelled effect is the thrust
augmentation acting through the contact arm --- a measured force scaled
by an established augmentation ratio --- and only $21\%$ of the gap is
a ground-\emph{induced moment}, which is the part that rests on the
rotor-interference model. The speculative claim is therefore a fifth of
the gap, not four-fifths of it.

One qualification belongs with the number. $\Delta M$ is not a
measurement of the aerodynamic moment: it is what the rigid balance
leaves over, and two of that balance's ingredients are themselves
modelled --- the collective thrust is $\sum C_T\omega_i^2$ from a bench
identification, hence an out-of-ground-effect value by construction,
and the circle fit of Sec.~V locates the \emph{kinematic} pivot,
whereas the balance needs the line of action of the normal-force
resultant. $\Delta M$ therefore collects the ground's total
contribution, aerodynamic and contact alike.

\paragraph{Two checks that use no external reference}
Both \eqref{eq:weight} and \eqref{eq:offset} close on themselves.
Eq.~\eqref{eq:offset} estimates the same offset twice, once from each
tip direction with that direction's own arm, and \eqref{eq:weight}
returns the weight from the moments. Neither closes without a
ground-effect term: the two directional offset estimates disagree by
$12.6$~mm in the median ($19.8$ at worst), and \eqref{eq:weight}
returns a weight $5.2\%$ below $mg$, low in all ten case/axis
configurations. Restoring the channels brings the directional spread
to $3.6$~mm ($9.0$ at worst) and the weight to $1.1\%$ below; the
thrust channel does most of the latter, since it enters
\eqref{eq:weight} as $(1+\eta_f^{GE})f_{\rm crit}$. Neither the load
cell nor the identified offsets enter these two checks.

The gap is also not an artifact of the onset detector. The four
model-agnostic baselines of Sec.~II-D bracket the identified moment
within a median $108$~mN\,m, and all detectors place the per-run onset
moment within $3$--$45$~mN\,m of one another, both well below
$\Delta M$.

\section*{The residual cannot be attributed, and three tests say why}

The $49$~mN\,m the model leaves is not a small coefficient error. Two
physically distinct explanations fit it equally well, and the data
cannot separate them.

\emph{H1, ground effect}: the residual is a mis-sized version of the
same two channels. \emph{H2, contact geometry}: the normal-force
resultant acts a distance $\delta l$ inboard of the kinematic pivot
circle, as a finite compliant contact patch implies, so the static
lever is $l_p - \delta l$. Both scale as force $\times$ length and are
degenerate in the constant part; the only discriminating structure is
the moment-proportional channel of H1, which grows with $|M_{\rm crit}|$
over a $4.3\times$ range across groups while H2 does not.
Table~\ref{tab:attrib} fits both.

\begin{table}[t]
\centering
\caption{Attribution of the threshold deficit over the twenty direction
groups. H1 free lets the data choose the channel gains; H1 fixed uses
the interference coefficients with a single scale.}
\label{tab:attrib}
\small
\setlength{\tabcolsep}{4pt}
\begin{tabular}{@{}l l r@{}}
\hline
hypothesis & fitted & residual \\
 & & RMS [mN\,m] \\
\hline
H1 free       & $\eta_f = 0.098\pm0.027$, $\eta_M = -0.055\pm0.063$ & $78.7$\\
H1 fixed      & scale $= 1.213\pm0.12$                              & $84.8$\\
H2            & $\delta l = 19.4\pm1.9$~mm                          & $83.6$\\
H1$+$H2       & const $= 0.235\pm0.06$~N\,m, $\eta_M = -0.031\pm0.052$ & $76.0$\\
\hline
\end{tabular}
\end{table}

Three things follow. Letting the gains free barely improves the fit,
$84.8$ to $78.7$~mN\,m, so the form is not the problem. The
discriminating gain is not detected at all: $\eta_M$ comes back
negative with an interval that spans zero, and the reason is visible in
the scale --- the $\eta_M M$ signal spans $18$ to $79$~mN\,m across
groups while the group-to-group scatter of the deficit is itself
$79$~mN\,m. And H2 alone fits as well as H1 does, $83.6$ against
$84.8$, with a lever displacement of $19.4\pm1.9$~mm.

That degeneracy is a property of the design, not an accident of the
sample, and three independent attempts to break it with the data in
hand all return null.

\emph{The rotation-centre height.} If a compliant leg moved the
instantaneous centre up from the foot, both the destabilising
coefficient and the pivot inertia would shrink together, which is what
the identified constants suggest. Freeing the pivot height in the
circle fit and sweeping it leaves the arc residual unchanged to two
decimals --- $0.14$~mm at the minimum against $0.14$ at ground level
--- because a $10^\circ$ arc trades the height against the radius
almost freely. The arm itself does bound it, however: a centre $25$~mm
above the ground would return $l_p = 154$~mm against the $140.4\pm3.6$
measured, so the centre is at the foot within about $10$~mm and this
reading is excluded rather than merely unresolved.

\emph{Ballast dependence.} Moving the ballast redistributes the load
among the contacts, so if the residual were contact compliance its
effective stiffness should track the offset. No component of the
offset predicts it: the strongest correlation over the ten
configurations is $|r| = 0.40$ at a permutation $p$ of $0.25$.

\emph{Contact migration during the arc.} The threshold is defined at
the onset, $\varphi = 0$, while the circle fit averages the arc that
follows, so a contact that migrates outboard as the vehicle tips would
inflate the fitted arm and over-predict the threshold --- the observed
sign. The migration leaves a trace, since the marker then traces a
trochoid rather than a circle, and the observed fit residuals of $0.1$
to $0.2$~mm correspond to an arm bias of $9.6$ to $18.4$~mm, which is
the right size against the $19.4$~mm fitted. But it does not vary the
right way: the correlation between residual and deficit is $+0.58$
pooled at $p = 0.081$, and it decomposes into $+0.75$ for roll against
$-0.67$ for pitch, so the pooled figure is the axis offset rather than
a relation.

A clearance sweep would separate H1 from H2, since the aerodynamic term
scales with $h/R$ and the contact term does not: repeating one axis at
$h = 0.60$~m would put the aerodynamic hypothesis at $0.069$~N\,m
against the contact hypothesis's unchanged $0.20$, a three-fold split
well above the $0.06$~N\,m floor. That sweep is outside the present
dataset, and Sec.~IX returns to it.

\section*{Why the deliverable does not need the attribution}

The reviewer's concern is whether this bias corrupts the delivered
offset. It does not, and the reason is structural rather than
fortunate.

\paragraph{The lever error cancels by direction}
Writing the two directional thresholds with their own arms and taking
the pivot-free average of \eqref{eq:pivotfree},
\begin{equation}
  M_{\rm ff} = \tfrac12(M_+ + M_-)
             = \tfrac12 (W-f)(l_+ - l_-) + S_{\rm off} ,
  \label{eq:ffcancel}
\end{equation}
so any error in the contact lever that is \emph{common} to the two tip
directions leaves $M_{\rm ff}$ untouched. The $19.4$~mm of H2 is such
an error, and it contributes nothing. Only the direction asymmetry
survives, and it is small: decomposing the residual of
Table~\ref{tab:attrib} by tip direction gives an antisymmetric part of
$54$~mN\,m in the median, which cancels, against a symmetric part of
$37$, which does not. The latter corresponds to $7.0$~mm of arm
asymmetry through $\mathrm{d}M/\mathrm{d}l = W - f = 10.6$~N, hence
$1.17$~mm of offset --- against the $1.43$~mm between-configuration
scatter of the offset error measured in Sec.~VIII-C.

That agreement is the point. What the static check fails to attribute
and what limits the delivered accuracy are the same quantity, seen from
two sides, and it is set by the contact geometry rather than by the
aerodynamics or by the estimator.

\paragraph{The geometry itself checks out}
The skid measures $286.29$~mm across and $250$~mm fore-aft, so the
half-widths are $143.14$ and $125.00$~mm against circle fits of
$140.5\pm3.6$ and $112.7\pm6.6$. Roll agrees to $2.7$~mm; pitch sits
$12.3$~mm inboard of the CAD edge, as a rail whose ends curve up
should. Substituting the CAD half-widths makes the deficit worse ---
pitch $69$ to $173$~mN\,m --- so the nominal arm is not the error, and
the per-run fitted arm is the right one to use. A rigid rolling foot
biases the fit in the right direction but by $0.4$ to $0.9$~mm for
realistic radii, an order of magnitude short of $\delta l$.

Worth recording separately is that the measured arm asymmetry
$l_+ - l_-$ is systematic per axis rather than scatter: negative in all
five roll groups and positive in all five pitch groups, $7.6$~mm in the
median. The pivot-free form of \eqref{eq:pivotfree} discards it by
assuming $l_+ = l_-$; the pivot-based form of \eqref{eq:offset} keeps
it, which is why the two are reported side by side below.

\paragraph{Effect on the identified offset}
Restoring both channels in the onset balance and pairing the
directions,
\begin{equation}
S_{\rm off} = \tfrac{1+\eta_M^{GE}}{2}(M_+ + M_-)
+ \tfrac{\eta_f^{GE}}{2}(f_+ l_+ - f_- l_-)
- \tfrac12\big[(W - f_+)l_+ - (W - f_-)l_-\big],
\label{eq:offge}
\end{equation}
the last bracket being dropped in the pivot-free variant.
Table~\ref{tab:offge} evaluates \eqref{eq:offge} against the load-cell
truth at both levels of physics and both treatments of the arms.

\begin{table}[t]
\centering
\caption{Identified CoM offset against the load-cell truth over the
ten case/axis configurations, with and without the ground-effect
channels of \eqref{eq:offge}. ``$\Delta$'' is the mean paired change
in $|$error$|$ against the $\eta^{GE} = 0$ row of the same block, with
the exact sign test beside it.}
\label{tab:offge}
\small
\setlength{\tabcolsep}{4pt}
\begin{tabular}{@{}ll rrr rr@{}}
\hline
arms & ground effect & RMS & mean & max & $\Delta$ & $p$ \\
 & & [mm] & [mm] & [mm] & [mm] & \\
\hline
pivot-free & none (reported) & $1.64$ & $-0.91$ & $3.34$ & --- & --- \\
pivot-free & single-rotor & $1.67$ & $-0.98$ & $3.44$ & $-0.00$ & $1.00$ \\
pivot-free & interference & $1.82$ & $-1.20$ & $3.77$ & $+0.19$ & $0.11$ \\
\hline
pivot-based & none & $1.77$ & $+0.42$ & $2.96$ & --- & --- \\
pivot-based & single-rotor & $1.73$ & $+0.35$ & $2.80$ & $-0.01$ & $1.00$ \\
pivot-based & interference & $1.65$ & $+0.13$ & $3.05$ & $-0.06$ & $1.00$ \\
\hline
\end{tabular}
\end{table}

The ground-effect term does not improve the offset, and no paired
comparison reaches significance: the correction is $0.41$~mm in the
median against a $0.45$~mm standard error of the comparison itself, so
this dataset bounds the ground-effect contribution to the offset at
about half a millimetre either way rather than resolving it. What does
move systematically is the mean bias, which the fullest inversion ---
arms and interference together --- removes, from $-0.91$ to
$+0.13$~mm at unchanged RMS.

\paragraph{What is applied, and what is reported}
The two uses of the same measurement need separating. The compensation
applied at take-off is $M_{\rm ff}$ of \eqref{eq:pivotfree} with no
ground-effect model anywhere: it is the in-situ balanced moment and
already carries whatever the ground contributes, which is precisely
what has to be cancelled, so subtracting a modelled term there would
remove part of the disturbance the compensation exists to reject.
Reading the same quantity as a mass property is the separate question
of Table~\ref{tab:offge}, and the offset is reported at $\eta^{GE} = 0$
to match the quantity flown, with the table standing as the bound on
what the omission costs.

The interval the excitation measures directly supports this. Any
moment strictly inside $[M_-, M_+]$ keeps both contacts loaded, the
midpoint maximises the smaller of the two margins, and both endpoints
come from the excitation rather than from a model. So $M_{\rm ff}$ is
the max--min-margin choice whatever the unmodelled terms turn out to
be, and by \eqref{eq:ffcancel} the direction-common part of them ---
$54$~mN\,m against a $1000$~mN\,m half-interval --- never enters it at
all.
\end{document}
